\section{Amenazas a la Validez}
\label{sec:amenazas-validez}

\subsection{Validez de constructo}
El p95 de latencia de k6 se usa como proxy de ``rendimiento percibido'', pero no se midió tiempo de
renderizado en el navegador ni Web Vitals bajo carga real de usuario (solo Lighthouse en condición de
reposo, sin carga concurrente). El puntaje SUS mide usabilidad percibida declarada, no eficiencia real de
tarea (tiempo para completar un flujo). Además, el resultado que se reporta proviene de la ronda del
18-sep ($n=15$, fecha sellada por un tercero, Sección~\ref{sec:evaluacion}); de la ronda anterior en
papel solo 4 de 15 hojas tenían fecha verificable y las otras 11 quedaron retractadas. Aun con la ronda
nueva, $n=15$ es una muestra pequeña, insuficiente para generalizar con confianza más allá de este
grupo, y sus participantes no son los mismos de la ronda en papel.

\subsection{Validez interna}
Las mediciones de rendimiento y Lighthouse se ejecutaron en la máquina de desarrollo personal de un
integrante del equipo. La hipótesis inicial (2026-08-17) fue contención de CPU por aplicaciones de
escritorio activas (Discord, Steam, Brave), confirmadas corriendo con \texttt{Get-Process | Sort CPU}.
Esa hipótesis se puso a prueba el 2026-08-30 en vez de repetirse sin verificar: se cerraron esas
aplicaciones y se remidió --- el promedio de Performance no cambió de forma apreciable, refutando la
hipótesis. El 2026-09-06 se identificó y verificó una causa real distinta para desktop: una imagen del
logo institucional de 333\,KB (2444$\times$2826\,px) mostrada en pantalla a 68\,px, y cabeceras de caché
de nginx que forzaban revalidación de red incluso para bundles con hash inmutables --- corregidas ambas,
Performance en desktop subió de 65 a 68.3. En mobile, sin embargo, el puntaje no se movió pese a que el
Largest Contentful Paint sí mejoró de forma real (7.6--7.8\,s $\to$ 6.7--6.8\,s): bajo el
\emph{throttling} de CPU 4$\times$ que aplica el perfil mobile de Lighthouse, el Total Blocking Time pasó
de 0\,ms a 90--100\,ms, evidencia de que el cuello de botella ahí es tiempo de ejecución de JavaScript del
framework (Angular + \texttt{zone.js}) sobre una CPU simulada más lenta, no peso de red --- una amenaza a
la validez interna con causa raíz ahora identificada y verificada (no solo descartada por eliminación,
como en la ronda anterior), pero que requiere un cambio de arquitectura (detección de cambios
\emph{zoneless}, \emph{server-side rendering}) fuera del alcance de esta entrega para cerrarse.

\subsection{Validez externa}
Los 50 usuarios concurrentes de k6 son un valor fijado por el RNF-01, no derivado de datos reales de uso
institucional de la UTEQ (que no existen para este sistema, al ser nuevo). No hay garantía de que ese
sea el nivel de concurrencia real durante un periodo de matrícula/defensas masivas. El usuario de
demostración y las pruebas manuales no reemplazan una prueba piloto con usuarios reales del dominio
(estudiantes, coordinadores) --- que es precisamente la brecha que el SUS busca cerrar, y que queda solo
parcialmente cerrada: las 15 respuestas de la ronda del 18-sep tienen fecha verificable, pero son
autoselección de quienes recibieron el enlace, sin muestreo ni control de perfil, y no cubren el ciclo
académico completo (Sección~\ref{sec:evaluacion}).

\subsection{Validez de conclusión}
El tamaño de muestra de Lighthouse ($n=3$ por perfil) es demasiado pequeño para una prueba de hipótesis
formal entre desktop y mobile (Sección~\ref{sec:evaluacion}); solo se reportan diferencias descriptivas.
La revisión de trabajos relacionados (Sección~\ref{sec:trabajos-relacionados}) no siguió un protocolo de
búsqueda sistemática formal con doble revisor, lo que limita la conclusión de que la cobertura de
literatura es exhaustiva --- se declaró explícitamente como una búsqueda dirigida, no una SLR completa.
Finalmente, el propio historial de datos fabricados detectados en este proyecto (Sección~\ref{sec:discusion},
RQ4) es una amenaza de conclusión de segundo orden: obliga a que cada cifra en este informe sea trazable
a \texttt{docs/mediciones/DATA-PROVENANCE.md}, precisamente porque la confianza por defecto en la
documentación de este proyecto específico ya se demostró insuficiente en al menos cinco episodios
distintos.
\newpage
